%%%%%%%%%%%%%%%%%%%%%%%%%%%%%%%%%%%%%%%%%%%%%%%%%%%%%%%%%%%%%%%%%%%%%%%%%%%%%%%%
% IEEE FG 2027 -- supplementary material (anonymised). Build: tectonic supplementary.tex
%%%%%%%%%%%%%%%%%%%%%%%%%%%%%%%%%%%%%%%%%%%%%%%%%%%%%%%%%%%%%%%%%%%%%%%%%%%%%%%%
\PassOptionsToPackage{table,dvipsnames}{xcolor}
\documentclass[a4paper, 10pt, conference, onecolumn]{ieeeconf}
\usepackage{FG2027}

%\FGfinalcopy % *** Uncomment this line for the final submission

\IEEEoverridecommandlockouts
\overrideIEEEmargins

% One column: the per-dataset tables and heatmaps are too wide for the
% two-column layout of the main paper.
\usepackage{iftex}
\ifPDFTeX\else\usepackage[T1]{fontenc}\fi
\usepackage{xcolor}
\usepackage{graphicx}
\usepackage{booktabs}
\usepackage{amsmath,amssymb}
\usepackage{multirow}
\usepackage{array}
\usepackage{adjustbox}
\usepackage[section]{placeins}
\usepackage[breaklinks,colorlinks,citecolor=blue,linkcolor=black,urlcolor=blue]{hyperref}

\newcommand{\tabtight}{\setlength{\tabcolsep}{3.5pt}\renewcommand{\arraystretch}{1.12}}
\newcommand{\TDS}{\textsc{TDS}}
\setlength{\emergencystretch}{4em}

\def\FGPaperID{****} % *** Enter the FG2027 Paper ID here

\title{\LARGE \bf
Supplementary Material\\[2pt]
\Large Measuring and Mitigating Spatiotemporal Aliasing in Video Action Recognition:\\
A Cross-Architecture Analysis at Scale
}

\author{Anonymous FG2027 submission\\ Paper ID \FGPaperID \\}

\begin{document}
\pagestyle{plain}
\maketitle

\renewcommand\thesection{S\arabic{section}}
\renewcommand\thefigure{S\arabic{figure}}
\renewcommand\thetable{S\arabic{table}}

This document provides supplementary material for the main paper. We include: (S1) full coverage$\times$stride heatmaps for all 8 models and 8 datasets; (S2) Levene variance inflation analysis; (S3) between-cell ANOVA effect sizes; (S4) action sensitivity taxonomy detail; (S5) confidence-cascade curves for all models; (S6) clip duration analysis; (S7) spectral correlation detail; (S8) implementation details (splits, sampling, frame normalization, target-resolution fine-tuning, positional embedding interpolation); (S9) full multi-dataset spatial resolution results (all 8 datasets $\times$ 8 models); (S10) interactive dashboard and code availability; (S11) \TDS{} ranking robustness: leave-one-architecture-out, architecture-family ablation, and bootstrap confidence intervals; (S12) spectral and label-space evidence-localization checks; (S13) additional robustness checks; (S14) frame selection and the under-sampled regime, with sampling intervals in seconds; (S15) the matched-evidence architecture comparison; (S16) the repeated-measures analysis of the coverage$\times$stride grid; (S17) held-out evaluation of the confidence cascade; (S18) convergence of low-resolution fine-tuning; (S19) sensitivity of the \TDS{} ranking to its definition; and (S20) multi-stride training as a mitigation.

% ── S1: Full heatmaps ─────────────────────────────────────────────────────────
\section{Full Coverage$\times$Stride Heatmaps}
\label{sec:sup_heatmaps}

Figures~\ref{fig:heatmap_r3d}--\ref{fig:heatmap_videomamba} show the complete coverage$\times$stride accuracy heatmaps for all 8 architectures across all 8 datasets. Each cell represents Top-1 accuracy (\%) at a specific (coverage, stride) configuration. Color scale: green=high accuracy, red=low accuracy. Datasets are ordered by decreasing \TDS{}.

Row order (coverage): 10\%, 25\%, 50\%, 75\%, 100\%. Column order (stride): 1, 2, 4, 8, 16.

\begin{figure}[!htbp]
  \centering
  \includegraphics[width=\textwidth]{images/sup_heatmap_r3d_18.pdf}
  \caption{R3D-18 (CNN-3D) coverage$\times$stride accuracy (\%) across 8 datasets. Sharp cliff at stride 4$\to$8 on AUTSL (75\%$\to$27\%) and similar cliff on FineGym are consistent with the high temporal demand of fine-grained actions. EPIC-Kitchens and UCF-101 retain high accuracy even at sparse sampling.}
  \label{fig:heatmap_r3d}
\end{figure}

\begin{figure}[!htbp]
  \centering
  \includegraphics[width=\textwidth]{images/sup_heatmap_mc3_18.pdf}
  \caption{MC3-18 (CNN-mix, channel-separated 3D) heatmaps. Similar pattern to R3D-18 but with slightly reduced aliasing on low-TDS datasets (UCF-101, DriveAct), consistent with the regularizing effect of 2D-channel separation. FineGym shows the same severe cliff pattern as AUTSL.}
  \label{fig:heatmap_mc3}
\end{figure}

\begin{figure}[!htbp]
  \centering
  \includegraphics[width=\textwidth]{images/sup_heatmap_r2plus1d_18.pdf}
  \caption{R2+1D-18 (separable spatiotemporal convolution) heatmaps. Higher base accuracy than R3D-18 on several datasets (SSv2: 42.6\% vs 37.1\%), but similar aliasing sensitivity pattern---cliff at stride 4$\to$8 on AUTSL and FineGym.}
  \label{fig:heatmap_r2p1d}
\end{figure}

\begin{figure}[!htbp]
  \centering
  \includegraphics[width=\textwidth]{images/sup_heatmap_slowfast_r50.pdf}
  \caption{SlowFast-R50 (dual-pathway) heatmaps. Cliff occurs at stride 2$\to$4 on AUTSL (77\%$\to$41\%)---one step earlier than other CNNs---because the fast pathway requires 32 frames. On FineGym, SlowFast achieves the highest base accuracy (78.2\%) among all models but suffers the worst aliasing loss (71.5pp at stride=16).}
  \label{fig:heatmap_slowfast}
\end{figure}

\begin{figure}[!htbp]
  \centering
  \includegraphics[width=\textwidth]{images/sup_heatmap_timesformer.pdf}
  \caption{TimeSformer (divided space--time attention) heatmaps. Notably flat across most stride values, with steep degradation only at the combined extreme of low coverage \emph{and} high stride. FineGym exception: drops 36.1pp at stride=16, revealing that even global divided attention cannot fully compensate for the 6--8Hz gymnastics phase transitions.}
  \label{fig:heatmap_timesformer}
\end{figure}

\begin{figure}[!htbp]
  \centering
  \includegraphics[width=\textwidth]{images/sup_heatmap_vivit.pdf}
  \caption{ViViT (factorized attention) heatmaps. Despite being a Transformer, ViViT aliases nearly as severely as 3D CNNs (AUTSL: cliff at stride 2$\to$4, identical to SlowFast). FineGym shows a 55.2pp drop at stride=16, consistent with factorized temporal attention being vulnerable to frame dropout in high-TDS domains.}
  \label{fig:heatmap_vivit}
\end{figure}

\begin{figure}[!htbp]
  \centering
  \includegraphics[width=\textwidth]{images/sup_heatmap_videomae.pdf}
  \caption{VideoMAE (masked autoencoder) heatmaps. High base accuracy (SSv2: 52.4\%, UCF-101: 95.4\%, FineGym: 78.0\%) but steep temporal degradation (AUTSL cliff at stride 4$\to$8; FineGym: 67.8pp drop). The MAE pretraining objective is spatial---reconstructing patches---and does not transfer to temporal robustness.}
  \label{fig:heatmap_videomae}
\end{figure}

\begin{figure}[!htbp]
  \centering
  \includegraphics[width=\textwidth]{images/sup_heatmap_videomamba.pdf}
  \caption{VideoMamba (SSM selective scan) heatmaps (all 8 datasets). On 5 of 8 datasets, VideoMamba shows a flat heatmap with minimal degradation across strides---the SSM recurrent state aggregates temporal context, providing natural robustness. FineGym (40.0pp), AUTSL (16.0pp), and SSv2 (14.1pp) are exceptions, revealing the limit of SSM temporal integration under high-frequency transitions.}
  \label{fig:heatmap_videomamba}
\end{figure}

% ── S2: Levene variance ───────────────────────────────────────────────────────
\section{Levene Variance Inflation Analysis}
\label{sec:sup_levene}

Figure~\ref{fig:levene} plots inter-class standard deviation at stride=1 versus stride=16 for all model-dataset pairs. Points above the diagonal indicate that sparse sampling increases inter-class variability; points below indicate that it decreases. Pairs with a significant Levene test ($p < 0.05$) are coloured by their variance ratio.

\begin{figure}[!htbp]
  \centering
  \includegraphics[width=0.8\columnwidth]{images/sup_levene_variance.pdf}
  \caption{Levene test: inter-class accuracy std at stride=1 vs.\ stride=16. Points above the dashed diagonal have higher variance under sparse sampling. Color encodes the variance ratio ($\sigma_{s16}/\sigma_{s1}$). The test is significant in 37 of 64 pairs (58\%; 22 pairs, 34\%, after Bonferroni correction). Variance inflates on HMDB-51 (6 pairs) and UCF-101 (3 pairs), most for VideoMAE/HMDB-51 (2.0$\times$) and R3D-18/HMDB-51 (1.85$\times$), and shrinks on SSv2, FineGym, AUTSL, Diving-48 and EPIC-Kitchens, where accuracy falls for nearly all classes. TimeSformer and VideoMamba stay within 0.9--1.14$\times$.}
  \label{fig:levene}
\end{figure}

Table~\ref{tab:levene_top} lists the 10 model-dataset pairs with the highest significant variance inflation, quantifying the instability masked by mean accuracy.

\begin{table}[!htbp]
\centering
\caption{The ten largest significant variance-inflation cases (stride=1$\to$16, coverage=100\%; Levene $p{<}0.05$).}
\label{tab:levene_top}
\scriptsize\tabtight
\begin{tabular}{l l r r r r}
\toprule
Model & Dataset & $\sigma_{s=1}$ & $\sigma_{s=16}$ & Ratio & Levene $p$ \\
\midrule
VideoMAE    & HMDB-51      & 0.125 & 0.253 & 2.02$\times$ & $<$0.0001 \\
R3D-18      & HMDB-51      & 0.164 & 0.303 & 1.85$\times$ & 0.0001 \\
ViViT       & HMDB-51      & 0.158 & 0.277 & 1.75$\times$ & 0.0004 \\
SlowFast    & HMDB-51      & 0.187 & 0.326 & 1.74$\times$ & $<$0.0001 \\
SlowFast    & UCF-101      & 0.146 & 0.240 & 1.65$\times$ & $<$0.0001 \\
R2+1D       & HMDB-51      & 0.205 & 0.310 & 1.51$\times$ & 0.0021 \\
MC3-18      & HMDB-51      & 0.177 & 0.262 & 1.48$\times$ & 0.0187 \\
ViViT       & UCF-101      & 0.110 & 0.147 & 1.34$\times$ & 0.0093 \\
VideoMAE    & AUTSL        & 0.172 & 0.229 & 1.33$\times$ & 0.0324 \\
VideoMAE    & UCF-101      & 0.094 & 0.121 & 1.29$\times$ & 0.0239 \\
\bottomrule
\end{tabular}
\end{table}

% ── S3: Full ANOVA tables ─────────────────────────────────────────────────────
\section{Between-Cell ANOVA Effect Size Tables}
\label{sec:sup_anova}

This section reports the classical between-cell decomposition, which treats the 25 grid cells of a model--dataset pair as independent observations and covers all 8 architectures. The main paper reports the repeated-measures analysis, detailed in Sec.~\ref{sec:sup_rmanova}; both order the architectures identically by stride effect. Figure~\ref{fig:anova_bar} shows the mean $\eta^2$ per model (averaged over all valid datasets), with error bars indicating standard deviation.

\begin{figure}[!htbp]
  \centering
  \includegraphics[width=0.9\columnwidth]{images/sup_anova_eta2.pdf}
  \caption{ANOVA effect sizes: $\eta^2(\text{coverage})$ + $\eta^2(\text{stride})$ per model. Coverage consistently dominates (blue bars, 50--88\% of variance). The stride effect is large for CNNs, ViViT and VideoMAE (red bars, 20--37\%) but small for VideoMamba and TimeSformer (9\%), matching the architecture-level aliasing differences in the main sweep. Error bars: std across all 8 datasets.}
  \label{fig:anova_bar}
\end{figure}

Table~\ref{tab:anova_full} provides the complete per-dataset $\eta^2$ values for all 8 datasets. All stride effects are significant ($p < 0.05$); coverage effects are significant ($p < 0.001$) in all cases. FineGym ($\eta^2_\text{stride}$ up to 0.46 for SlowFast) and AUTSL show the largest stride effects of all 8 datasets, consistent with their co-equal top ranking on TDS (FineGym=55.9pp, AUTSL=53.0pp; see main paper Sec.~IV-A and supplementary~S11 for the leave-one-out sensitivity of their relative order).

\begin{table}[!htbp]
\centering
\caption{Between-cell ANOVA $\eta^2$ per model and dataset. Cov = $\eta^2(\text{coverage})$, Str = $\eta^2(\text{stride})$. ``---'': not available for that pair.}
\label{tab:anova_full}
\scriptsize\tabtight
\begin{adjustbox}{max width=\textwidth}
\begin{tabular}{l *{8}{cc}}
\toprule
\multirow{2}{*}{Model} & \multicolumn{2}{c}{AUTSL} & \multicolumn{2}{c}{FineGym} & \multicolumn{2}{c}{SSv2} & \multicolumn{2}{c}{Diving-48} & \multicolumn{2}{c}{HMDB-51} & \multicolumn{2}{c}{DriveAct} & \multicolumn{2}{c}{EPIC}
 & \multicolumn{2}{c}{UCF-101} \\
\cmidrule(lr){2-3}\cmidrule(lr){4-5}\cmidrule(lr){6-7}\cmidrule(lr){8-9} \cmidrule(lr){10-11}\cmidrule(lr){12-13}\cmidrule(lr){14-15}\cmidrule(lr){16-17} & Cov & Str & Cov & Str & Cov & Str & Cov & Str & Cov & Str & Cov & Str & Cov & Str
 & Cov & Str \\
\midrule
VMamba & --- & --- & .76 & .17 & .88 & .09 & .94 & .04 & .90 & .09 & .83 & .15 & .95 & .05 & .88 & .07 \\
TSF    & .91 & .05 & .73 & .20 & .87 & .10 & .97 & .02 & .92 & .07 & .84 & .14 & .91 & .08 & .85 & .07 \\
MC3    & .66 & .20 & .48 & .34 & .54 & .28 & .81 & .12 & .80 & .18 & .75 & .22 & .89 & .11 & .78 & .17 \\
R3D    & .60 & .21 & .46 & .32 & .60 & .26 & .77 & .15 & .68 & .30 & .58 & .33 & .75 & .24 & .70 & .25 \\
R2+1D  & .63 & .19 & .46 & .34 & .62 & .26 & .76 & .17 & .72 & .26 & .49 & .38 & .79 & .18 & .76 & .20 \\
VMAE   & .69 & .18 & .51 & .29 & .69 & .22 & .82 & .12 & .84 & .16 & .43 & .43 & .87 & .11 & .81 & .14 \\
ViViT  & .61 & .24 & .49 & .35 & .51 & .32 & .69 & .20 & .71 & .27 & .60 & .31 & .80 & .18 & .68 & .26 \\
SF     & .51 & .29 & .30 & .46 & .44 & .36 & .51 & .33 & .55 & .39 & .39 & .46 & .67 & .29 & .61 & .36 \\
\bottomrule
\end{tabular}
\end{adjustbox}
\end{table}

% ── S4: Taxonomy ──────────────────────────────────────────────────────────────
\section{Action Sensitivity Taxonomy}
\label{sec:sup_taxonomy}

\begin{figure}[!htbp]
  \centering
  \includegraphics[width=\textwidth]{images/sup_taxonomy.pdf}
  \caption{Aliasing loss (pp) per sensitivity tier and dataset. Each bar shows mean accuracy drop (stride=1$\to$16) within the tier, annotated with the number of action classes. UCF-101 Low tier: $-0.3$pp (static, appearance-dominated actions lose essentially no accuracy under sparse sampling). AUTSL: all three tiers remain high ($>38$pp), consistent with sign language being high-demand across action subclasses.}
  \label{fig:taxonomy_full}
\end{figure}

Table~\ref{tab:taxonomy_detail} lists the threshold values that define tier boundaries per dataset, derived from the 33rd and 67th percentiles of the per-class accuracy drop.

\begin{table}[!htbp]
\centering
\caption{Taxonomy tier thresholds (accuracy drop, pp): classes below the Low threshold alias negligibly; classes above the High threshold are severely affected. Thresholds are dataset-specific; per-class drops are averaged over the 8 architectures before the percentiles are taken.}
\label{tab:taxonomy_detail}
\scriptsize\tabtight
\begin{tabular}{l r r r r r}
\toprule
Dataset & Classes & Low threshold & High threshold & High count & Low count \\
\midrule
FineGym       &  97 & 50.4pp & 71.2pp & 33 & 32 \\
AUTSL         & 226 & 46.7pp & 57.3pp & 75 & 74 \\
SSv2          & 174 & 20.6pp & 31.8pp & 58 & 57 \\
DriveAct      &  33 & 14.0pp & 29.6pp & 11 & 11 \\
Diving-48     &  47 & 12.6pp & 22.4pp & 16 & 16 \\
HMDB-51       &  51 &  7.9pp & 20.2pp & 17 & 17 \\
EPIC-Kitchens &  89 &  0.0pp &  7.1pp & 30 & 11 \\
UCF-101       & 101 &  1.2pp &  6.2pp & 35 & 31 \\
\bottomrule
\end{tabular}
\end{table}

% ── S5: Routing curves ────────────────────────────────────────────────────────
\section{Confidence-Cascade Curves --- All Models}
\label{sec:sup_routing}

Figure~\ref{fig:routing_all} shows accuracy vs.\ average frames on SSv2 for the confidence cascade, which thresholds the maximum class probability of the stride-4 pass. Each curve is the result of sweeping confidence threshold $\tau \in [0.10, 0.95]$. Dashed horizontal lines indicate fixed-budget baselines (4f, 16f) and the oracle.

\begin{figure}[!htbp]
  \centering
  \includegraphics[width=\textwidth]{images/sup_cascade_all_models.pdf}
  \caption{Confidence cascade: accuracy vs.\ avg frames (SSv2, stride-4 cheap $\to$ stride-1 dense). For architectures with flat 4f$\approx$16f accuracy (TimeSformer, VideoMamba), routing near-trivially selects 4f for most videos and matches dense accuracy. For architectures with large 4f--16f gaps (SlowFast, ViViT), the routing curve provides meaningful operating points between the fixed-budget extremes.}
  \label{fig:routing_all}
\end{figure}

Table~\ref{tab:routing_full} provides the cascade results across all models and datasets at the threshold that is optimal on the evaluated split (avg frames $\leq 8$); the held-out counterpart is in Sec.~\ref{sec:sup_routing_heldout}.

\begin{table}[!htbp]
\centering
\caption{Confidence cascade with in-sample threshold selection: best accuracy at avg $\leq8$ frames per model/dataset. Values are Top-1 accuracy (\%). Bold = highest per column. ``---'': not available for that pair.}
\label{tab:routing_full}
\scriptsize\tabtight
\begin{adjustbox}{max width=\textwidth}
\begin{tabular}{l cccccccc r}
\toprule
Model & AUTSL & FineGym & Diving & SSv2 & HMDB & DriveAct & EPIC & UCF & Avg \\
\midrule
R3D-18       & 73.3 & 60.2 & 27.6 & 34.6 & 80.2 & 66.1 & 36.7 & 81.6 & 57.5 \\
MC3-18       & 63.4 & 62.7 & 31.8 & 32.4 & 78.5 & \textbf{69.4} & 36.5 & 85.6 & 57.5 \\
R2+1D        & 74.2 & 65.3 & 36.3 & 40.8 & 73.3 & 59.6 & 34.8 & 89.0 & 59.2 \\
SlowFast     & 59.5 & 45.2 & 46.6 & 35.0 & 76.3 & 67.4 & 37.4 & 87.9 & 56.9 \\
TimeSformer  & 67.1 & 65.4 & 38.9 & 42.5 & 80.6 & 68.5 & 32.4 & 91.3 & 60.8 \\
ViViT        & 67.2 & 54.4 & 48.9 & 32.2 & 80.1 & 64.3 & 31.3 & 94.3 & 59.1 \\
VideoMAE     & \textbf{79.5} & 67.4 & \textbf{52.6} & \textbf{51.9} & \textbf{84.7} & 69.0 & \textbf{37.8} & \textbf{95.9} & \textbf{67.4} \\
VideoMamba   & --- & \textbf{72.4} & 37.4 & 44.4 & 69.8 & 57.6 & 28.4 & 88.5 & 56.9 \\
\bottomrule
\end{tabular}
\end{adjustbox}
\end{table}

% ── S6: Clip duration ─────────────────────────────────────────────────────────
\section{Clip Duration Analysis}
\label{sec:sup_duration}

\begin{figure}[!htbp]
  \centering
  \includegraphics[width=\textwidth]{images/sup_clip_duration.pdf}
  \caption{Aliasing loss (pp) vs.\ clip duration bin for all 8 architectures across multiple datasets. Durations use frame rates measured from the video files. Counter-intuitive finding: \textbf{shorter clips alias more} (Pearson $r=-0.3$ to $-0.8$ across models). Short clips have less temporal redundancy; any frame drop removes a larger fraction of the available evidence. Long clips contain natural temporal repetition that compensates for sparse sampling.}
  \label{fig:duration}
\end{figure}

% ── S7: Spectral correlation ──────────────────────────────────────────────────
\section{Spectral Correlation Detail}
\label{sec:sup_spectral}

We estimated per-class temporal frequency demand using optical-flow magnitude (Farneback algorithm~\cite{farneback2003two}) computed on 5 representative videos per class, sampling 16 frames per video at 112$\times$112px. The mean inter-frame flow magnitude serves as a proxy for temporal activity level. This analysis is exploratory and should not be read as confirmatory evidence for the Nyquist-aliasing hypothesis on its own; see the main paper (Sec.~IV-A, Fig.~3b) for the summary figure including FineGym.

Table~\ref{tab:spectral} reports Pearson correlation between per-class mean flow magnitude and aliasing loss (accuracy drop stride=1$\to$16) per dataset. Prior studies on small closed-vocabulary datasets (UCF-101, Kinetics-400) have reported near-perfect optical-flow--aliasing correlations ($r{=}0.99$) using controlled synthetic flow protocols. Our large-scale evaluation on 8 diverse datasets yields a more nuanced picture ($r=-0.475$ to $0.327$). This has at least four plausible explanations: (1) complex background motion in egocentric datasets (EPIC-Kitchens) decorrelates flow from action content; (2) sign language (AUTSL) involves localized, subtle movements that optical flow at 112px partially misses; (3) our use of raw flow magnitude rather than class-conditioned frequency analysis is a coarser proxy; (4) FineGym's negative correlation suggests that for discrete-event classes (a held position, a precise landing), raw motion magnitude is the wrong signal entirely---what matters is whether sparse sampling catches the one brief moment that defines the class, which low-flow (slow, precise) actions are more likely to lose than high-flow (fast, sustained) ones.

\begin{table}[!htbp]
\centering
\caption{Spectral correlation: Pearson $r$ between optical-flow magnitude and per-class aliasing loss. Positive $r$ supports the Nyquist hypothesis; near-zero or negative values suggest flow magnitude is an insufficient (or inverted) proxy for temporal demand.}
\label{tab:spectral}
\scriptsize\tabtight
\begin{tabular}{l r r r}
\toprule
Dataset & Pearson $r$ & $p$-value & Significant ($p{<}0.05$) \\
\midrule
DriveAct      & 0.327 & 0.063 & Marginal \\
AUTSL         & 0.285 & $<$0.001 & Yes \\
SSv2          & 0.184 & 0.015 & Yes \\
HMDB-51       & 0.180 & 0.207 & No \\
Diving-48     & 0.140 & 0.346 & No \\
UCF-101       & 0.031 & 0.762 & No \\
EPIC-Kitchens & --0.002 & 0.988 & No \\
FineGym       & --0.475 & $<$0.0001 & Yes (inverted) \\
\bottomrule
\end{tabular}
\end{table}

% ── S8: Implementation details ────────────────────────────────────────────────
\section{Implementation Details}
\label{sec:sup_impl}

\paragraph{Hardware}
Training experiments used NVIDIA A100 (40GB) and H200 NVL (141GB) GPUs (CUDA 13.2). Temporal-sweep evaluation (1,600 configurations) was parallelized across 8 nodes using a custom Slurm auto-submitter. Inference latency measurements used an NVIDIA RTX PRO 6000 Blackwell GPU (96GB) with bfloat16 autocast, batch size=1, 20 warmup iterations, 100 benchmark iterations, and CUDA event timing with explicit synchronization.

\paragraph{Splits, clip sampling, and evaluation determinism}
We use the official train/validation/test split when a dataset provides one. When the public split is larger than needed for the controlled sweep, we form a deterministic balanced evaluation subset with approximately 20 validation clips per class; if a class has fewer available clips, all available validation clips are used. The same clip list is reused across all architectures, resolutions, coverages, and strides for a dataset. This makes pairwise stride/resolution differences comparable and avoids model-specific sampling noise. Reported main-paper numbers are validation-set measurements; no test-set tuning is used.

\paragraph{Seeds and uncertainty}
Because the full grid requires thousands of model--dataset--resolution evaluations, each model--dataset fine-tuning run uses one fixed training seed and deterministic evaluation. Confidence intervals in S11 are therefore not random-seed intervals: they bootstrap over the 8-architecture pool. Per-class taxonomy and Levene analyses use class-level variation, and cascade thresholds are calibrated and scored on disjoint halves of the validation split (Sec.~\ref{sec:sup_routing_heldout}).

\paragraph{Frame and crop normalization}
Training uses random resized crop, horizontal flip, and ImageNet normalization. Evaluation uses deterministic resize/center-crop at the requested spatial resolution. For every coverage/stride setting, frames are selected from the same leading candidate window. Each architecture receives its expected input length; if a candidate window is too short, we pad by repeating the last available frame, so stride changes alter sampling density rather than tensor shape.

\paragraph{Training hyperparameters}
All models fine-tuned for 10 epochs from Kinetics-400 checkpoints. Transformers and VideoMamba: AdamW, lr=$2\times10^{-5}$, weight decay=0.05, batch size=32 per GPU. CNNs: AdamW, lr=$10^{-4}$, batch size=128 per GPU. SlowFast: AdamW, lr=$10^{-4}$, batch size=32 per GPU. All models use cosine learning rate decay with linear warmup (1 epoch). Spatial augmentation: random resized crop, horizontal flip, normalization (ImageNet statistics).

\paragraph{Positional embedding interpolation for Transformers and SSMs}
For target-resolution fine-tuning at non-native resolutions, Transformer and SSM positional embeddings must be adapted. \textit{The wrong approach}: using the HuggingFace mismatched-size flag discards positional embeddings entirely, causing a large accuracy cliff (typically 10--20pp worse than our approach). \textit{Our approach}: load the native-resolution checkpoint, then bicubically interpolate spatial positional embeddings from the native grid ($14{\times}14{=}196$ tokens at 224px) to the target grid ($\lfloor r/16\rfloor^2$ tokens at resolution $r$) before copying weights to a model instantiated at the target resolution. This preserves the spatial topology learned during pretraining and ensures fine-tuning starts from a meaningful positional encoding rather than a random one. The same bicubic interpolation applies to TimeSformer, VideoMAE, ViViT, and VideoMamba.

\paragraph{SlowFast spatial pooling fix}
The original SlowFast-R50 head uses a fixed $7\times7$ average pooling kernel, assuming 224px input. For non-native resolutions, we replace both head pooling layers with \texttt{AdaptiveAvgPool3d(1)}, enabling forward passes at arbitrary input resolutions without parameter changes.

\paragraph{FineGym dataset details}
FineGym~\cite{shao2020finegym} contains 97 fine-grained gymnastics action classes (floor exercise, uneven bars, balance beam, vault; parallel bars, horizontal bar, pommel horse, rings for men). We use the FineGym288 subset (1,960 clips), with deterministic balanced clip lists for the controlled sweeps. Preprocessing: clips are trimmed to the annotated action boundary (mean duration 3.2s at 25fps) and padded with the last frame if shorter than 16 frames.

% ── S9: Multi-dataset spatial resolution ──────────────────────────────────────
\section{Multi-Dataset Spatial Resolution Results}
\label{sec:sup_spatial}

Tables~\ref{tab:sup_sp_ssv2}--\ref{tab:sup_sp_finegym} extend the spatial robustness analysis from the main paper (Sec.~IV-E) to all 8 datasets. Tables \ref{tab:sup_sp_ssv2}--\ref{tab:sup_sp_epic} show spatial-resolution results at 96--224px. Table~\ref{tab:sup_sp_finegym} shows FineGym target-resolution fine-tuning results at 48--224px (one fine-tuned checkpoint per resolution). \textbf{Bold} = native training resolution. $\Delta_\text{max}$ = largest accuracy drop from native across all tested resolutions.

\begin{table}[!htbp]
\centering
\caption{Spatial robustness on \textbf{SSv2} (repeated from main paper; no retraining).}
\label{tab:sup_sp_ssv2}
\scriptsize\tabtight
\begin{adjustbox}{max width=\textwidth}
\begin{tabular}{l r rrrr r}
\toprule
Model & Native & 96px & 112px & 160px & 224px & $\Delta_\text{max}$ \\
\midrule
R3D-18      & 112px & 35.9 & \textbf{37.1} & 30.3 & 17.2 & $-19.9$pp \\
MC3-18      & 112px & 32.4 & \textbf{33.6} & 27.8 & 14.8 & $-18.8$pp \\
R2+1D       & 112px & 40.1 & \textbf{42.6} & 36.2 & 20.8 & $-21.8$pp \\
SlowFast    & 224px & 27.7 & 32.4 & 45.7 & \textbf{49.5} & $-21.8$pp \\
\midrule
TimeSformer & 224px & 39.3 & 39.4 & 41.4 & \textbf{42.3} & $-3.0$pp \\
ViViT       & 224px & 36.4 & 37.1 & 38.1 & \textbf{38.3} & $-1.9$pp \\
VideoMAE    & 224px & 48.9 & 49.2 & 51.5 & \textbf{52.3} & $-3.4$pp \\
VideoMamba  & 224px & 37.5 & 39.3 & 42.5 & \textbf{43.9} & $-6.4$pp \\
\bottomrule
\end{tabular}
\end{adjustbox}
\end{table}

\begin{table}[!htbp]
\centering
\caption{Spatial robustness on \textbf{UCF-101} (no retraining).}
\label{tab:sup_sp_ucf}
\scriptsize\tabtight
\begin{adjustbox}{max width=\textwidth}
\begin{tabular}{l r rrrr r}
\toprule
Model & Native & 96px & 112px & 160px & 224px & $\Delta_\text{max}$ \\
\midrule
R3D-18      & 112px & 79.9 & \textbf{81.2} & 75.3 & 55.1 & $-26.1$pp \\
MC3-18      & 112px & 82.9 & \textbf{85.4} & 81.0 & 63.9 & $-21.5$pp \\
R2+1D       & 112px & 88.0 & \textbf{88.6} & 85.1 & 64.9 & $-23.7$pp \\
SlowFast    & 224px & 69.7 & 75.8 & 87.5 & \textbf{91.4} & $-21.7$pp \\
\midrule
TimeSformer & 224px & 89.1 & 89.7 & 90.7 & \textbf{91.5} & $-2.4$pp \\
ViViT       & 224px & 92.0 & 93.5 & 94.7 & \textbf{95.1} & $-3.1$pp \\
VideoMAE    & 224px & 95.3 & 95.7 & 96.1 & \textbf{96.4} & $-1.1$pp \\
VideoMamba  & 224px & 76.7 & 78.7 & 82.6 & \textbf{83.5} & $-6.8$pp \\
\bottomrule
\end{tabular}
\end{adjustbox}
\end{table}

\begin{table}[!htbp]
\centering
\caption{Spatial robustness on \textbf{HMDB-51} (no retraining).}
\label{tab:sup_sp_hmdb}
\scriptsize\tabtight
\begin{adjustbox}{max width=\textwidth}
\begin{tabular}{l r rrrr r}
\toprule
Model & Native & 96px & 112px & 160px & 224px & $\Delta_\text{max}$ \\
\midrule
R3D-18      & 112px & 78.8 & \textbf{80.3} & 75.9 & 59.8 & $-20.5$pp \\
MC3-18      & 112px & 76.7 & \textbf{78.6} & 74.2 & 51.0 & $-27.6$pp \\
R2+1D       & 112px & 74.4 & \textbf{73.1} & 68.4 & 50.5 & $-22.6$pp \\
SlowFast    & 224px & 59.5 & 67.2 & 80.5 & \textbf{81.6} & $-22.1$pp \\
\midrule
TimeSformer & 224px & 75.9 & 77.2 & 77.5 & \textbf{78.9} & $-3.0$pp \\
ViViT       & 224px & 78.0 & 79.0 & 79.9 & \textbf{79.5} & $-1.5$pp \\
VideoMAE    & 224px & 83.7 & 83.1 & 84.0 & \textbf{83.8} & $-0.7$pp \\
VideoMamba  & 224px & 59.9 & 63.8 & 67.2 & \textbf{69.8} & $-9.9$pp \\
\bottomrule
\end{tabular}
\end{adjustbox}
\end{table}

\begin{table}[!htbp]
\centering
\caption{Spatial robustness on \textbf{Diving-48} (no retraining).}
\label{tab:sup_sp_div}
\scriptsize\tabtight
\begin{adjustbox}{max width=\textwidth}
\begin{tabular}{l r rrrr r}
\toprule
Model & Native & 96px & 112px & 160px & 224px & $\Delta_\text{max}$ \\
\midrule
R3D-18      & 112px & 26.5 & \textbf{28.8} & 22.1 & 12.8 & $-16.0$pp \\
MC3-18      & 112px & 30.8 & \textbf{31.6} & 27.3 & 13.7 & $-17.9$pp \\
R2+1D       & 112px & 33.8 & \textbf{35.3} & 30.1 & 18.7 & $-16.6$pp \\
SlowFast    & 224px & 18.4 & 25.0 & 43.3 & \textbf{50.5} & $-32.1$pp \\
\midrule
TimeSformer & 224px & 37.4 & 39.3 & 38.6 & \textbf{38.1} & $-0.7$pp \\
ViViT       & 224px & 48.0 & 49.9 & 52.8 & \textbf{53.0} & $-5.0$pp \\
VideoMAE    & 224px & 48.7 & 48.3 & 50.9 & \textbf{48.6} & $-0.3$pp \\
VideoMamba  & 224px & 29.2 & 29.1 & 36.1 & \textbf{36.6} & $-7.5$pp \\
\bottomrule
\end{tabular}
\end{adjustbox}
\end{table}

\begin{table}[!htbp]
\centering
\caption{Spatial resolution results on \textbf{AUTSL} (sign language). CNN degradation is severe under off-native evaluation; SlowFast drops from 79.2\% at 224px to 22.5\% at 96px ($-56.7$pp).}
\label{tab:sup_sp_autsl}
\scriptsize\tabtight
\begin{adjustbox}{max width=\textwidth}
\begin{tabular}{l r rrrr r}
\toprule
Model & Native & 96px & 112px & 160px & 224px & $\Delta_\text{max}$ \\
\midrule
R3D-18      & 112px & 72.2 & \textbf{75.0} & 63.0 & 27.9 & $-47.1$pp \\
MC3-18      & 112px & 59.1 & \textbf{63.7} & 52.7 & 12.2 & $-51.5$pp \\
R2+1D       & 112px & 73.7 & \textbf{75.9} & 68.9 & 38.4 & $-37.5$pp \\
SlowFast    & 224px & 22.5 & 32.6 & 70.0 & \textbf{79.2} & $-56.7$pp \\
\midrule
TimeSformer & 224px & 65.1 & 66.0 & 68.3 & \textbf{68.5} & $-3.4$pp \\
ViViT       & 224px & 49.2 & 51.8 & 53.3 & \textbf{53.6} & $-4.4$pp \\
VideoMAE    & 224px & 77.3 & 78.2 & 79.5 & \textbf{79.6} & $-2.3$pp \\
VideoMamba  & 224px & 30.4 & 34.2 & 53.3 & \textbf{65.8} & $-35.4$pp \\
\bottomrule
\end{tabular}
\end{adjustbox}
\end{table}

\begin{table}[!htbp]
\centering
\caption{Spatial robustness on \textbf{DriveAct} (in-cabin activities; no retraining).}
\label{tab:sup_sp_drive}
\scriptsize\tabtight
\begin{adjustbox}{max width=\textwidth}
\begin{tabular}{l r rrrr r}
\toprule
Model & Native & 96px & 112px & 160px & 224px & $\Delta_\text{max}$ \\
\midrule
R3D-18      & 112px & 65.0 & \textbf{68.3} & 52.7 & 19.0 & $-49.3$pp \\
MC3-18      & 112px & 65.2 & \textbf{69.0} & 38.4 &  6.2 & $-62.8$pp \\
R2+1D       & 112px & 67.4 & \textbf{62.5} & 40.0 & 12.9 & $-49.6$pp \\
SlowFast    & 224px & 31.7 & 40.2 & 65.8 & \textbf{70.1} & $-38.4$pp \\
\midrule
TimeSformer & 224px & 68.3 & 69.2 & 69.6 & \textbf{70.3} & $-2.0$pp \\
ViViT       & 224px & 59.4 & 58.7 & 62.1 & \textbf{63.4} & $-4.7$pp \\
VideoMAE    & 224px & 71.4 & 69.9 & 71.7 & \textbf{73.0} & $-3.1$pp \\
VideoMamba  & 224px & 35.0 & 38.2 & 46.9 & \textbf{57.8} & $-22.8$pp \\
\bottomrule
\end{tabular}
\end{adjustbox}
\end{table}

\begin{table}[!htbp]
\centering
\caption{Spatial robustness on \textbf{EPIC-Kitchens} (egocentric cooking; no retraining).}
\label{tab:sup_sp_epic}
\scriptsize\tabtight
\begin{adjustbox}{max width=\textwidth}
\begin{tabular}{l r rrrr r}
\toprule
Model & Native & 96px & 112px & 160px & 224px & $\Delta_\text{max}$ \\
\midrule
R3D-18      & 112px & 33.3 & \textbf{37.2} & 26.7 &  8.6 & $-28.6$pp \\
MC3-18      & 112px & 36.1 & \textbf{36.2} & 27.2 &  9.4 & $-26.8$pp \\
R2+1D       & 112px & 32.8 & \textbf{35.5} & 27.6 & 11.8 & $-23.7$pp \\
SlowFast    & 224px & 11.3 & 17.7 & 31.3 & \textbf{37.5} & $-26.2$pp \\
\midrule
TimeSformer & 224px & 20.5 & 19.2 & 25.6 & \textbf{31.7} & $-12.5$pp \\
ViViT       & 224px & 30.2 & 34.1 & 36.1 & \textbf{37.9} & $-7.7$pp \\
VideoMAE    & 224px & 30.9 & 32.6 & 36.5 & \textbf{38.6} & $-7.7$pp \\
VideoMamba  & 224px & 11.1 & 12.6 & 20.8 & \textbf{28.3} & $-17.2$pp \\
\bottomrule
\end{tabular}
\end{adjustbox}
\end{table}

\begin{table}[!htbp]
\centering
\caption{\textbf{FineGym} spatial resolution results after target-resolution fine-tuning (one checkpoint per resolution, 10 fine-tuning epochs from native checkpoint). Bold = native training resolution (112px for CNNs, 224px for Transformers/SSM). Unlike other datasets in S9 (cross-resolution without retraining), FineGym shows CNN performance that \emph{improves beyond native} after fine-tuning at higher resolutions (R3D-18: 71.5\%@112px$\to$76.1\%@224px), consistent with higher spatial resolution aiding gymnastics pose discrimination.}
\label{tab:sup_sp_finegym}
\scriptsize\tabtight
\begin{adjustbox}{max width=\textwidth}
\begin{tabular}{l r rrrrr}
\toprule
Model & Native & 48px & 96px & 112px & 160px & 224px \\
\midrule
R3D-18      & 112px & 56.4 & 72.6 & \textbf{71.5} & 74.9 & 76.1 \\
MC3-18      & 112px & 46.3 & 62.3 & \textbf{64.3} & 73.8 & 65.7 \\
R2+1D       & 112px & 61.6 & 73.1 & \textbf{76.8} & 79.4 & 76.2 \\
SlowFast    & 224px & 36.3 & 67.2 & 69.6 & 76.3 & \textbf{78.2} \\
\midrule
TimeSformer & 224px & 24.3 & 48.9 & 50.8 & 62.2 & \textbf{66.4} \\
ViViT       & 224px & 21.3 & 39.1 & 40.8 & 49.6 & \textbf{69.9} \\
VideoMAE    & 224px & 25.4 & 51.1 & 57.5 & 64.8 & \textbf{78.0} \\
VideoMamba  & 224px & 29.1 & 51.6 & 58.2 & 66.0 & \textbf{73.2} \\
\bottomrule
\end{tabular}
\end{adjustbox}
\end{table}

% ── S10: Dashboard and code availability ──────────────────────────────────────
\section{Interactive Dashboard and Code Availability}
\label{sec:sup_dashboard}

\paragraph{Interactive visualization dashboard}
An interactive dashboard accompanies this work; its public URL will be released after the review period. The dashboard provides:
\begin{enumerate}
\item \textbf{Temporal explorer}: coverage$\times$stride heatmaps per model and
  dataset, with configurable scaling and stride annotations.
\item \textbf{Spatial browser}: accuracy curves across resolutions for all
  8 models and datasets; target resolution tuning by family.
\item \textbf{Recommender}: given a target dataset, latency budget, frame rate,
  and deployment context, the system applies deterministic rules from the paper
  tables to recommend a model--resolution--stride operating point. Any optional
  natural-language summarizer is disabled by default for anonymous review and is
  not used for any scientific result.
\item \textbf{TDS comparison}: rank any subset of the 8 datasets by temporal demand,
  with model-specific breakdowns and Spearman correlation display.
\item \textbf{Routing simulator}: confidence threshold sweep for any model on any dataset,
  with visual Pareto-frontier display.
\end{enumerate}


\paragraph{Code availability}
Evaluation scripts, model fine-tuning code, result tables, and the dashboard source, including the target-resolution pipeline with bicubic positional embedding interpolation, will be released in a public repository after review.

% ── S11: Leave-one-architecture-out TDS stability ─────────────────────────────
\section{Leave-One-Architecture-Out \texorpdfstring{\TDS{}}{TDS} Stability}
\label{sec:sup_loo}

To assess whether the \TDS{} ranking depends on the specific 8-architecture pool evaluated, we recompute \TDS{} for all 8 datasets after removing each architecture in turn (7-architecture pools) and compare the resulting ranking to the full 8-architecture ranking via Spearman $\rho$.

\begin{table}[!htbp]
\centering
\caption{Leave-one-architecture-out \TDS{} ranking stability. $\rho$ is the Spearman correlation between the 8-dataset ranking computed with the named architecture excluded (7-architecture pool) and the full 8-architecture ranking.}
\label{tab:loo_tds}
\scriptsize\tabtight
\begin{tabular}{l r r}
\toprule
Excluded architecture & Spearman $\rho$ & $p$ \\
\midrule
R3D-18       & 1.000 & $<10^{-4}$ \\
MC3-18       & 1.000 & $<10^{-4}$ \\
R2+1D        & 1.000 & $<10^{-4}$ \\
SlowFast     & 1.000 & $<10^{-4}$ \\
TimeSformer  & 1.000 & $<10^{-4}$ \\
ViViT        & 1.000 & $<10^{-4}$ \\
VideoMAE     & 1.000 & $<10^{-4}$ \\
VideoMamba   & 0.976 & $3.3\times10^{-5}$ \\
\bottomrule
\end{tabular}
\end{table}

For 7 of 8 leave-one-out pools, the ranking is \emph{exactly} identical to the full-pool ranking ($\rho=1.00$). Removing VideoMamba is the single exception ($\rho=0.976$): it swaps the \#1/\#2 positions, since excluding VideoMamba raises AUTSL's average (58.3pp, from 53.0pp) and lowers FineGym's average (58.2pp, from 55.9pp)---VideoMamba is comparatively robust on AUTSL (16.0pp, below its own median drop) and comparatively fragile on FineGym (40.8pp, its worst dataset), so its presence in the pool is precisely what tips FineGym narrowly ahead of AUTSL. Ranks 3--8 (SSv2 through UCF-101) are unaffected by any single exclusion. We therefore report FineGym and AUTSL as co-equal highest-\TDS{} domains in the main paper rather than asserting a strict order, while the rest of the ranking is robust under leave-one-architecture-out by this test.

\subsection{Architecture-Family Ablation}
\label{sec:sup_family_ablation}

Leave-one-out removes a single architecture at a time; a stronger test removes an entire family. We recompute \TDS{} under three restricted pools: \textbf{CNN-only} ($n{=}4$: R3D-18, MC3-18, R2+1D, SlowFast), \textbf{Transformer-only} ($n{=}3$: TimeSformer, ViViT, VideoMAE), and \textbf{Transformer+SSM} ($n{=}4$: adding VideoMamba). Table~\ref{tab:family_ablation} reports \TDS{} under each pool and Spearman $\rho$ against the full ranking.

\begin{table}[!htbp]
\centering
\caption{\TDS{} per dataset under architecture-family ablation pools, compared to the full 8-architecture ranking (sorted by full-pool \TDS{}).}
\label{tab:family_ablation}
\scriptsize\tabtight
\begin{tabular}{l r r r r}
\toprule
Dataset & Full ($n{=}8$) & CNN-only ($n{=}4$) & Transf.-only ($n{=}3$) & Transf.+SSM ($n{=}4$) \\
\midrule
FineGym       & 55.92 & 62.05 & 53.06 & 49.80 \\
AUTSL         & 53.02 & 67.20 & 46.45 & 38.83 \\
SSv2          & 27.29 & 32.16 & 25.95 & 22.43 \\
DriveAct      & 21.48 & 24.05 & 23.21 & 18.92 \\
Diving-48     & 19.16 & 21.74 & 20.59 & 16.59 \\
HMDB-51       & 16.47 & 21.47 & 14.54 & 11.47 \\
EPIC-Kitchens &  9.77 & 12.89 &  8.27 &  6.64 \\
UCF-101       &  4.86 &  7.36 &  2.90 &  2.35 \\
\midrule
Spearman $\rho$ vs.\ full & --- & 0.976 & 1.000 & 1.000 \\
\bottomrule
\end{tabular}
\end{table}

The ranking is essentially invariant to architecture-family composition: removing all four CNNs, keeping only the three Transformers, or keeping Transformers+SSM all leave ranks 3--8 in identical order, and only the CNN-only pool reproduces the same FineGym/AUTSL swap already characterized in the leave-one-out test above ($\rho=0.976$ in both cases, for the same reason: CNNs are comparatively more fragile on AUTSL's fine handshape transitions than on FineGym's phase boundaries, which narrows the gap). This directly addresses whether \TDS{} reflects a property of the dataset pool versus an artifact of including any one architecture family: under every family-restricted pool we tested, the ranking is the same dataset ordering reported in the main paper, with the sole exception of the already-disclosed \#1/\#2 ambiguity.

\subsection{Bootstrap Confidence Intervals}
\label{sec:sup_bootstrap}

To quantify ranking uncertainty directly rather than only testing discrete leave-one-out/family pools, we bootstrap over architectures: we resample 8 architectures with replacement ($B{=}10{,}000$ resamples) from the original pool, recompute \TDS{} for each dataset on each resample, and report the resulting 95\% percentile confidence interval in Table~\ref{tab:bootstrap_ci}.

\begin{table}[!htbp]
\centering
\caption{Bootstrap 95\% CI for \TDS{} per dataset (resampling 8 architectures with replacement, $B{=}10{,}000$).}
\label{tab:bootstrap_ci}
\scriptsize\tabtight
\begin{tabular}{l r r}
\toprule
Dataset & \TDS{} (bootstrap mean) & 95\% CI \\
\midrule
FineGym       & 55.99 & [47.24, 63.86] \\
AUTSL         & 53.13 & [36.02, 67.03] \\
SSv2          & 27.35 & [20.40, 33.92] \\
DriveAct      & 21.57 & [13.92, 29.35] \\
Diving-48     & 19.23 & [10.67, 28.37] \\
HMDB-51       & 16.54 & [9.53, 24.21] \\
EPIC-Kitchens &  9.80 & [6.25, 13.22] \\
UCF-101       &  4.88 & [2.18, 8.36] \\
\bottomrule
\end{tabular}
\end{table}

All eight 95\% CIs are disjoint from zero, indicating non-zero temporal demand at the architecture-pool level. AUTSL's CI is the widest (width 31.0pp) because its per-architecture drops are the most heterogeneous (TimeSformer and VideoMamba near-flat, CNNs collapse sharply), while EPIC-Kitchens and UCF-101 have the narrowest CIs, consistent with their low and relatively architecture-independent \TDS{}.

We additionally bootstrap the FineGym$-$AUTSL gap directly (same resampling procedure, computing $\TDS(\text{FineGym})-\TDS(\text{AUTSL})$ on each resample): mean $=+2.92$pp, 95\% CI $=[-4.58,\,+11.73]$, with FineGym exceeding AUTSL in 73.9\% of resamples. The confidence interval includes zero, which is the appropriate statistical test of the co-equal-ranking claim: the data are consistent with FineGym narrowly leading AUTSL in expectation, but do not support rejecting the null hypothesis of equal \TDS{} at the 95\% level. This is the quantitative basis, rather than only the leave-one-out swap, for reporting FineGym and AUTSL as co-equal highest-\TDS{} domains in the main paper.

% ── S12: Spectral and evidence-localization checks ───────────────────────────
\section{Spectral and Evidence-Localization Checks}
\label{sec:sup_nyquist}

The main paper's Nyquist--Shannon framing is conceptual: we characterize aliasing through accuracy degradation under controlled stride and coverage sweeps rather than by directly measuring action-frequency spectra. Here we report two diagnostics. The first is a deliberately literal whole-frame optical-flow spectral proxy, which fails in the opposite direction from a naive signal-level Nyquist prediction. The second asks whether classifier evidence is temporally localized, which is better aligned with \TDS{} but should still be read as a label-space diagnostic, not classical spectral validation.

\paragraph{Methodology} For each of 8 datasets and 5 resolutions (48, 96, 112, 160, 224\,px), we sample up to 12 classes $\times$ 2 videos. We decode up to 64 \emph{consecutive} frames in native order, resize to the target resolution, and compute inter-frame optical-flow magnitude (Farneback) as a temporal signal sampled once per original frame. We estimate each signal's spectral cutoff frequency $f_c$ (90\% cumulative periodogram power, cycles/frame) and average over sampled videos to obtain one $f_c$ per (dataset, resolution) pair. We pair each $f_c$ with empirical stride-sensitivity at the same resolution: mean accuracy drop from stride=1 to stride=16 at coverage=100\%, averaged over 8 architectures and read from the existing sweep. This yields $8\times5=40$ (dataset, resolution) points.

\begin{table}[!htbp]
\centering
\caption{Mean flow-derived spectral cutoff frequency $f_c$, videos per resolution, and mean empirical stride-sensitivity, averaged over the 5 tested resolutions per dataset and sorted by stride-sensitivity.}
\label{tab:nyquist_v2}
\scriptsize\tabtight
\begin{tabular}{l r r r}
\toprule
Dataset & Mean $f_c$ & $n$/res. & Mean stride-sens. \\
\midrule
FineGym       & 0.188 & 24 & 46.30 \\
AUTSL         & 0.121 & 24 & 45.61 \\
SSv2          & 0.339 & 24 & 25.70 \\
DriveAct      & 0.192 & 21 & 18.65 \\
HMDB-51       & 0.383 & 24 & 16.57 \\
Diving-48     & 0.182 & 24 & 15.21 \\
EPIC-Kitchens & 0.248 & 23 &  9.70 \\
UCF-101       & 0.387 & 24 &  5.45 \\
\bottomrule
\end{tabular}
\end{table}

\paragraph{Result} Pooling across all 40 (dataset, resolution) points, $f_c$ is significantly \emph{negatively} correlated with stride-sensitivity (Spearman $\rho=-0.635$, permutation $p=1.0\times10^{-5}$, $n=40$)---the opposite sign from the naive prediction that higher whole-frame temporal-frequency content should require denser sampling. FineGym integrates cleanly into this pattern: it has the highest mean stride-sensitivity (46.30pp) but a moderate whole-frame cutoff frequency ($f_c=0.188$ cycles/frame), close to DriveAct and Diving-48 rather than UCF-101/HMDB-51.

\paragraph{Interpretation} This is the same qualitative inversion already reported for FineGym in the main paper ($r=-0.475$), now shown across all 8 datasets with real statistical power rather than being a single-dataset anomaly. UCF-101 and HMDB-51---appearance- and sports-driven domains with large whole-body and camera motion---have the highest flow-derived cutoff frequencies but among the lowest stride-sensitivity, because the action remains identifiable from almost any sampled frame. AUTSL and FineGym show the opposite behavior: their discriminative evidence is localized in handshape changes or brief gymnastics phase events, which are diluted by a whole-frame flow aggregate. We conclude that whole-frame optical-flow frequency content is a systematically \emph{anti-correlated}, not merely noisy, proxy for the temporal demand that matters to a classifier.

\subsection*{Label-Space Evidence-Localization Checks}

The inversion above arises because whole-frame optical flow measures \emph{pixel-level motion magnitude}, not the temporal localisation of \emph{discriminative evidence}. We conduct three complementary tests on FineGym ($n=97$ classes) using the fine-tuned TimeSformer checkpoint, which probe whether the model's evidence is temporally concentrated.

\paragraph{Test A: Temporal burstiness of optical flow (CPU)} Per-class, we compute the temporal coefficient of variation of inter-frame flow magnitudes ($\mathrm{CV} = \sigma_\mathrm{time}/\mu_\mathrm{time}$). After controlling for mean flow via partial correlation, the temporal standard deviation positively predicts aliasing: $\partial r = +0.235$, $p=0.021$. This suggests that burstiness, not mean flow magnitude, is closer to the classifier-evidence signal, though the confound ($r(\sigma,\mu)=+0.909$) limits interpretability.

\paragraph{Test B: TimeSformer temporal attention entropy (GPU)} We register forward hooks on the temporal self-attention Dropout layer in all 12 TimeSformer blocks, capturing the $[196\times12\times8\times8]$ temporal attention weight tensor per clip. The mean entropy of the time-over-time attention distribution (Nat, over spatial patches and heads) positively correlates with aliasing: $r=+0.317$, $p=0.0015$. Classes where the model spreads temporal attention broadly alias more, consistent with the idea that evidence distributed across several brief instants is more vulnerable to frame removal.

\paragraph{Test C: Sliding-window confidence concentration (GPU, most direct)} For each clip we decode 48 consecutive frames, slide a window of width $T=8$ (the model input length) with stride 4, yielding 11 windows per clip, and record the softmax probability assigned to the correct class at each position. The temporal \emph{concentration} of this confidence profile ($1 - H/\log n_\mathrm{windows}$) measures how localised in time the model finds class evidence. Concentration positively predicts aliasing: $r=+0.263$, $p=0.009$. After controlling for mean optical-flow magnitude, the partial correlation strengthens: $\partial r = +0.385$, $p<0.001$ ($n=97$ classes).

\paragraph{Summary} Table~\ref{tab:label_space_nyquist} collects the three results.

\begin{table}[!htbp]
\centering
\caption{Label-space evidence-localization checks on FineGym ($n=97$ classes). All three tests use a fine-tuned TimeSformer trained on FineGym. Partial $r$ controls for whole-frame optical-flow magnitude.}
\label{tab:label_space_nyquist}
\scriptsize\tabtight
\begin{tabular}{l l r r r r}
\toprule
Test & Feature & $r$ & $p$ & partial $r$ & $p$ \\
\midrule
A: Optical flow burstiness & $\sigma_t(\text{flow}) / \mu_t(\text{flow})$ & $-0.421$ & ${<}0.001$ & $+0.235$ & $0.021$ \\
B: Temporal attention entropy & Mean attention $H$ (nats) & $+0.317$ & $0.0015$ & --- & --- \\
C: Sliding-window confidence & Confidence concentration & $+0.263$ & $0.009$ & $\mathbf{+0.385}$ & $\mathbf{<0.001}$ \\
\bottomrule
\end{tabular}
\end{table}

Test C provides the clearest classifier-level diagnostic: temporal concentration of \emph{correct-class confidence} in a fine-tuned model positively predicts per-class aliasing after removing the motion-level confound. This supports an evidence-localization interpretation: classes where the model is confident only within a narrow temporal window are more likely to lose that window under coarse stride. It does not establish literal spectral aliasing of the video signal.

\subsection*{Generalisation of Test C to All 8 Datasets}
\label{sec:sup_sw_generalisation}

We generalise the sliding-window confidence-concentration test (Test C) from the single FineGym/TimeSformer case to all 8 datasets and an ensemble of architectures. For each (model, dataset) we decode 48 consecutive frames per clip, slide a window of width $T$ (the model's input length) with stride chosen to yield ${\approx}11$ windows, and record the correct-class softmax probability at each position; per class we average the concentration $1-H/\log n_\mathrm{windows}$ over clips. We then correlate per-class concentration with the per-class aliasing loss (mean accuracy drop stride=1$\to$16 from the taxonomy) via Spearman $\rho$.

\paragraph{Window-geometry constraint} The metric requires the decoded sequence to be substantially longer than the window $T$, so that windows are well separated and concentration reflects genuine temporal localisation. With 48 decoded frames, models with $T\le16$ satisfy this (neighbouring-window overlap $\le81\%$), but the two models with $T=32$ (ViViT, SlowFast) collapse to stride 1 and $97\%$ overlap: their 17 windows differ by a single frame, the confidence profile is near-flat, and the concentration measure floors out (mean concentration $0.10$ vs.\ $0.22$--$0.36$ for $T\le16$ models). On short-clip datasets such as FineGym (median ${\approx}37$ source frames) a 32-frame window cannot be slid at all. We therefore restrict the ensemble to the 6 architectures with $T\le16$ (TimeSformer, VideoMAE, R3D-18, MC3-18, R2+1D, VideoMamba); the $T=32$ models are a measurement artifact of the protocol rather than a property of the data.

\paragraph{Result} Table~\ref{tab:sw_ensemble} reports the ensemble correlation per dataset. The relationship is \emph{positive in all 8 datasets} (pooled Spearman $\rho=+0.257$), and significant ($p<0.05$) in 4: HMDB-51, UCF-101, DriveAct, and SSv2. The four non-significant datasets are positive but weaker: FineGym ($\rho=+0.14$; Pearson $+0.20$, $p=0.045$) is damped by its visually homogeneous gymnastics domain, EPIC-Kitchens by egocentric camera motion, AUTSL because aliasing there is governed by spatial gesture trajectories rather than temporal localisation, and Diving-48 because all classes share the same approach--flight--entry structure, leaving little inter-class variance in concentration. We therefore treat this as converging evidence that temporal concentration of recognizer evidence is related to aliasing susceptibility, not as a causal or sampling-theoretic proof.

\begin{table}[!htbp]
\centering
\caption{Sliding-window confidence-concentration test generalised to all 8 datasets. Spearman $\rho$ between per-class confidence concentration and per-class aliasing loss, using a 6-architecture ensemble with window $T\le16$ (ViViT and SlowFast, $T=32$, excluded; see window-geometry constraint). Positive in 8/8 datasets; \textbf{bold} = $p<0.05$. SSv2 uses 4 models because the sliding-window pass for VideoMAE and R3D-18 on SSv2 was unavailable.}
\label{tab:sw_ensemble}
\scriptsize\tabtight
\begin{tabular}{l r r r r}
\toprule
Dataset & Classes & Models & Spearman $\rho$ & $p$ \\
\midrule
HMDB-51       &  51 & 6 & \textbf{+0.606} & ${<}0.001$ \\
UCF-101       & 101 & 6 & \textbf{+0.442} & ${<}0.001$ \\
DriveAct      &  33 & 6 & \textbf{+0.430} & 0.013 \\
SSv2          & 174 & 4 & \textbf{+0.269} & ${<}0.001$ \\
FineGym       &  97 & 6 & +0.138 & 0.177 \\
EPIC-Kitchens &  89 & 6 & +0.093 & 0.386 \\
AUTSL         & 226 & 6 & +0.074 & 0.270 \\
Diving-48     &  47 & 6 & +0.001 & 0.994 \\
\midrule
Pooled mean   &     &   & \textbf{+0.257} & --- \\
\bottomrule
\end{tabular}
\end{table}

Generalisation of Test B (TimeSformer temporal attention entropy) to the remaining datasets requires per-dataset attention extraction and is left to future work; the sliding-window concentration test (Test C), being architecture-agnostic, is the one we generalise here.

% ── S13: Additional robustness checks ───────────────────────────────────────
\section{Additional Robustness Checks}
\label{sec:sup_extra_checks}

\paragraph{Observational architecture interpretation}
The architecture comparison in the main paper is intentionally phrased as an empirical association rather than a controlled causal ablation. Attention layout, state-space recurrence, tokenization, pretraining objective, parameter count, input length, native resolution, and optimization all covary across the evaluated architectures. A definitive mechanism test would require holding the backbone fixed while swapping only temporal aggregation modules. Our contribution is the large controlled sampling sweep showing that, within this architecture pool, models with global divided attention or recurrent state integration are consistently less sensitive to temporal subsampling than CNNs and factorized-attention Transformers.

\paragraph{Dense-accuracy headroom does not explain \TDS{}}
Because several datasets are evaluated on balanced subsets or fixed 20-clips-per-class protocols, absolute gaps to external leaderboard numbers are not directly comparable. The relevant risk is whether under-optimized model--dataset pairs inflate \TDS{} simply because they have more room to lose accuracy. We therefore compute an internal dense-accuracy headroom: for each dataset, the gap between each model's stride=1 accuracy and the best stride=1 accuracy achieved by any model in our pool. If under-optimization drove aliasing loss, this headroom should correlate positively with stride=1$\to$16 drop. It does not: across the 64 model--dataset pairs, Spearman $\rho=-0.22$ ($p=0.079$); at the dataset level, mean headroom vs.\ \TDS{} gives $\rho=0.36$ ($p=0.385$).

\begin{table}[!htbp]
\centering
\caption{Dense-accuracy headroom check. Headroom is the mean gap from each model's stride=1 accuracy to the best stride=1 accuracy achieved by any model on the same dataset. It is not significantly correlated with \TDS{}, suggesting that temporal sensitivity is not an artifact of under-optimized dense accuracy.}
\label{tab:headroom_check}
\scriptsize\tabtight
\begin{tabular}{l r r r r}
\toprule
Dataset & Mean stride=1 & Best stride=1 & Mean headroom & \TDS{} \\
\midrule
FineGym       & 73.8 & 79.9 &  6.1 & 55.9 \\
AUTSL         & 72.9 & 82.3 &  9.3 & 53.0 \\
SSv2          & 42.5 & 52.4 &  9.9 & 27.6 \\
DriveAct      & 67.4 & 73.9 &  6.5 & 21.9 \\
Diving-48     & 40.3 & 53.0 & 12.7 & 19.2 \\
HMDB-51       & 78.2 & 84.0 &  5.8 & 16.6 \\
EPIC-Kitchens & 34.9 & 39.4 &  4.5 &  9.7 \\
UCF-101       & 89.0 & 95.5 &  6.5 &  4.9 \\
\bottomrule
\end{tabular}
\end{table}

\paragraph{Coverage--stride interaction}
Figures in S1 visualize the complete coverage$\times$stride grids. To test additivity, we average accuracy over architectures for one high-\TDS{} and one low-\TDS{} dataset, then decompose each 5$\times$5 grid into row/column effects plus an interaction residual. FineGym has a meaningful interaction component: 15.0\% of grid variance, with RMSE 8.3pp. Low coverage is poor regardless of stride, while high coverage exposes the stride cliff. UCF-101 is closer to additive (4.0\%; RMSE 1.7pp).

\begin{table}[!htbp]
\centering
\caption{Mean accuracy (\%) over architectures for a high-\TDS{} dataset (FineGym) and a low-\TDS{} dataset (UCF-101). Rows are coverage; columns are stride.}
\label{tab:coverage_stride_examples}
\scriptsize\tabtight
\begin{adjustbox}{max width=\textwidth}
\begin{tabular}{l rrrrr @{\hskip18pt} l rrrrr}
\toprule
\multicolumn{6}{c}{FineGym} & \multicolumn{6}{c}{UCF-101} \\
\cmidrule(lr){1-6}\cmidrule(lr){7-12}
Coverage & s1 & s2 & s4 & s8 & s16 & Coverage & s1 & s2 & s4 & s8 & s16 \\
\midrule
10\%  &  8.0 &  7.2 &  6.4 &  5.8 &  5.8 & 10\%  & 74.5 & 73.2 & 69.5 & 64.2 & 57.7 \\
25\%  & 16.9 & 15.6 & 12.4 &  8.3 &  6.4 & 25\%  & 81.7 & 81.6 & 80.1 & 75.0 & 67.8 \\
50\%  & 42.9 & 40.7 & 28.7 & 17.1 &  9.3 & 50\%  & 86.6 & 86.6 & 85.7 & 82.9 & 76.6 \\
75\%  & 63.4 & 59.6 & 42.3 & 25.1 & 13.1 & 75\%  & 88.4 & 88.7 & 88.5 & 86.8 & 81.4 \\
100\% & 73.8 & 68.0 & 49.8 & 29.8 & 15.7 & 100\% & 89.0 & 89.0 & 89.0 & 88.1 & 84.1 \\
\bottomrule
\end{tabular}
\end{adjustbox}
\end{table}

\paragraph{Confidence-threshold calibration}
The cascade is not trained; it reuses two fixed inference outputs (stride-4 cheap and stride-1 dense) and a scalar confidence threshold. All cascade numbers in the main paper select the threshold on a calibration half of the validation split and score it on the other half (Sec.~\ref{sec:sup_routing_heldout}); the curves in Sec.~\ref{sec:sup_routing} show the full threshold sweep for reference.

% ── S14: Frame selection ─────────────────────────────────────────────────────
\section{Frame Selection and the Under-Sampled Regime}
\label{sec:sup_padding}

This section details how frames are selected for each (coverage, stride)
setting and quantifies the regime in which stride alters the model input.

\paragraph{Sampler} For a clip of $T$ frames at coverage $\mathcal{C}$ and
stride $s$, with model frame budget $B$, the candidate pool is
$\{0,s,2s,\dots\}$ within the first $\lfloor T\mathcal{C}/100\rfloor$ frames.
If the pool holds at least $B$ entries it is resampled uniformly to exactly $B$;
otherwise the shortfall is filled by repeating the last candidate.
Table~\ref{tab:sup_indices} illustrates both cases. Frame budgets are 8 for
TimeSformer and VideoMamba, 16 for R3D-18, MC3-18, R(2+1)D and VideoMAE, and 32
for SlowFast and ViViT.

\begin{table}[!htbp]
\centering
\caption{Frame positions supplied to a $B{=}16$ model at $\mathcal{C}{=}100\%$.
For a long clip the selected positions are nearly independent of stride; for a
short clip a large stride leaves fewer than $B$ distinct frames.}
\label{tab:sup_indices}
\scriptsize\tabtight
\begin{tabular}{l l r l}
\toprule
Clip & $s$ & Distinct & Positions (last six) \\
\midrule
$T{=}400$ & 1  & 16 & 266, 293, 319, 346, 372, 399 \\
$T{=}400$ & 4  & 16 & 264, 292, 316, 344, 368, 396 \\
$T{=}400$ & 16 & 16 & 256, 288, 304, 336, 352, 384 \\
\midrule
$T{=}60$  & 1  & 16 & 39, 43, 47, 51, 55, 59 \\
$T{=}60$  & 8  &  8 & 56, 56, 56, 56, 56, 56 \\
$T{=}60$  & 16 &  4 & 48, 48, 48, 48, 48, 48 \\
\bottomrule
\end{tabular}
\end{table}

\paragraph{Sampling interval in seconds} Table~\ref{tab:sup_seconds} expresses
the stride axis physically. Frame rates are measured from the video files. At a
fixed nominal stride the realised interval differs by up to $3.9\times$ across
datasets, so the stride at which a dataset becomes under-sampled depends on its
frame rate and clip length as well as on its content.

\begin{table}[!htbp]
\centering
\caption{Median seconds between distinct sampled frames at
$\mathcal{C}{=}100\%$ for a 16-frame model.}
\label{tab:sup_seconds}
\scriptsize\tabtight
\begin{tabular}{l r rrrrr}
\toprule
Dataset & fps & $s{=}1$ & $s{=}2$ & $s{=}4$ & $s{=}8$ & $s{=}16$ \\
\midrule
UCF-101       & 25 & 0.395 & 0.395 & 0.395 & 0.395 & 0.617 \\
Diving-48     & 25 & 0.265 & 0.265 & 0.265 & 0.314 & 0.600 \\
SSv2          & 12 & 0.255 & 0.255 & 0.325 & 0.620 & 1.167 \\
DriveAct      & 15 & 0.183 & 0.183 & 0.256 & 0.500 & 1.000 \\
HMDB-51       & 30 & 0.167 & 0.167 & 0.167 & 0.260 & 0.504 \\
AUTSL         & 30 & 0.125 & 0.125 & 0.131 & 0.252 & 0.475 \\
EPIC-Kitchens & 50 & 0.120 & 0.120 & 0.120 & 0.158 & 0.301 \\
\bottomrule
\end{tabular}
\end{table}

\paragraph{Where the stride effect arises} Restricting to clips whose candidate
pool still holds at least $B$ frames at $s{=}16$, the stride $1\to16$ accuracy
difference is $-0.5$pp on average, against $+10.7$pp over all clips of the same
model--dataset pairs (21 pairs with at least 20 such clips). The accuracy loss
measured by the sweep is therefore the loss incurred when the stride no longer
leaves $B$ distinct frames in the observed window. On AUTSL, SSv2 and DriveAct
every clip is in this regime at $s{=}16$; on UCF-101, 75.5\% of clips are.

\paragraph{Relation to \TDS{}} Across datasets, the share of clips in the
under-sampled regime at $s{=}16$ correlates with \TDS{} at Spearman $0.96$. The
relationship is gating rather than explanatory: among AUTSL, SSv2 and DriveAct,
all of which are entirely in the regime, \TDS{} still spans $21.9$--$53.0$pp.
Clip length relative to frame budget determines whether temporal demand is
expressed; dataset content determines how large it is.

% ── S15: Matched-evidence protocol ───────────────────────────────────────────
\section{Matched-Evidence Architecture Comparison}
\label{sec:sup_matched}

Frame budget $B$ determines the stride at which each architecture becomes
under-sampled, and $B$ is rank-correlated with the mean stride drop of the main
paper's temporal table (Spearman $0.849$, $p{=}0.008$). To compare architectures
at equal evidence, we supply every architecture the same $k$ distinct frames at
the same positions, uniformly spanning the clip, resampled to each model's own
input length by uniform index interpolation rather than end-repetition. Only
$k\leq8$ is strictly matched, since the smallest budget in the pool is 8. The
comparison covers the seven datasets other than FineGym.

\begin{table}[!htbp]
\centering
\caption{Top-1 accuracy (\%) at $k$ matched distinct frames, averaged over seven
datasets. ``Rel.\ loss'' is the accuracy lost from $k{=}8$ to $k{=}2$ relative
to $k{=}8$; ``Stride drop'' is the mean stride $1\to16$ drop of the main paper.
Rows are sorted by stride drop.}
\label{tab:sup_matched}
\scriptsize\tabtight
\begin{tabular}{l r rrrr r r}
\toprule
Model & $B$ & $k{=}1$ & $k{=}2$ & $k{=}4$ & $k{=}8$ & Rel.\ loss & Stride drop \\
\midrule
TimeSformer &  8 & 21.1 & \textbf{33.8} & \textbf{52.1} & \textbf{58.6} & 42.3\% & 10.3pp \\
VideoMamba  &  8 & 21.9 & 28.4 & 41.0 & 44.5 & 36.2\% & 11.2pp \\
MC3-18      & 16 & \textbf{24.7} & 29.5 & 39.8 & 51.0 & 42.2\% & 22.6pp \\
R3D-18      & 16 & 18.2 & 24.8 & 36.6 & 50.1 & 50.6\% & 29.7pp \\
R(2+1)D     & 16 & 18.0 & 25.3 & 37.5 & 50.5 & 50.0\% & 30.0pp \\
ViViT       & 32 & 20.4 & 27.6 & 41.3 & 53.6 & 48.4\% & 30.0pp \\
VideoMAE    & 16 & 17.5 & 23.2 & 38.2 & 48.7 & 52.3\% & 32.8pp \\
SlowFast    & 32 & 18.4 & 24.1 & 32.6 & 38.5 & 37.3\% & 42.1pp \\
\bottomrule
\end{tabular}
\end{table}

Table~\ref{tab:sup_matched} gives the result. TimeSformer is the most accurate
architecture at every $k\geq2$, by $4.3$, $10.8$ and $5.0$pp over the runner-up,
and SlowFast is the least accurate at $k\geq4$. Accuracy at matched evidence is
positively rank-correlated with the stride-drop ordering at every $k$: Spearman
$0.64$ ($k{=}1$), $0.86$ ($k{=}2$, $p{=}0.007$), $0.62$ ($k{=}4$) and $0.48$
($k{=}8$). With eight architectures only the $k{=}2$ value is individually
significant, so we read these as a consistent direction rather than as four
independent confirmations.

The accuracy lost between $k{=}8$ and $k{=}2$ is not associated with $B$
(Spearman $0.23$, $p{=}0.58$), nor does it reproduce the stride-drop ordering
($0.21$, $p{=}0.61$): relative losses lie between 36\% and 52\% for all eight
architectures. What separates architectures under matched evidence is therefore
the accuracy they deliver per frame rather than the fraction they lose.
VideoMamba ranks third up to $k{=}4$ but seventh at $k{=}8$, which indicates
that its small stride drop owes part to its 8-frame input.

Two caveats apply. Architectures with a large $B$ receive far fewer distinct
frames than they were trained on, which places them further from their design
point than the 8-frame models. And the comparison remains observational with
respect to temporal aggregation design, as discussed in
Sec.~\ref{sec:sup_extra_checks}. The table is reproduced by
\texttt{scripts/fg2027/matched\_evidence\_table.py}.

% ── S16: Repeated-measures ANOVA ─────────────────────────────────────────────
\section{Repeated-Measures Analysis of the Coverage\texorpdfstring{$\times$}{x}Stride Grid}
\label{sec:sup_rmanova}

All 25 cells of the grid score the same clips. We therefore analyse the grid
with the clip as subject and coverage and stride as fully crossed within-subject
factors, testing each effect against its own subject$\times$factor error term
and reporting partial $\eta^2$. Only clips present in all 25 cells are used, so
the design is balanced. The analysis covers the seven architectures whose
per-clip outputs are available (56 model--dataset pairs); VideoMamba is reported
in the between-cell analysis of Sec.~\ref{sec:sup_anova}.

Partial $\eta^2$ is $0.292\pm0.152$ for coverage, $0.232\pm0.145$ for stride and
$0.070\pm0.081$ for their interaction. The interaction is largest on the
high-demand datasets, reaching $0.271$ on SlowFast/AUTSL, $0.262$ on
R3D-18/AUTSL and $0.232$ on SlowFast/FineGym. The per-model ordering of the
stride effect (TimeSformer $0.124$, MC3-18 $0.197$, R3D-18 $0.226$, R(2+1)D
$0.238$, ViViT $0.247$, VideoMAE $0.259$, SlowFast $0.331$) matches the
between-cell analysis, so the grouping of architectures does not depend on the
choice of error term.

% ── S17: Held-out cascade ────────────────────────────────────────────────────
\section{Held-Out Evaluation of the Confidence Cascade}
\label{sec:sup_routing_heldout}

The cascade returns the stride-4 prediction when its maximum class probability
exceeds $\tau$ and the stride-1 prediction otherwise. Cost is counted in
acquired frames (16 for stride 1, 4 for stride 4). We choose $\tau$ on a
stratified calibration half of the validation split, subject to an average cost
of at most 8 frames, and score it on the other half, over $1000$ resamples.
Intervals are 95\% percentile intervals over resamples. Seven architectures have
the per-clip outputs this requires.

\begin{table}[!htbp]
\centering
\caption{Held-out cascade against fixed dense inference. Left: mean over the 8
datasets per architecture, with the number of datasets whose interval includes
zero. Right: TimeSformer per dataset.}
\label{tab:sup_cascade}
\scriptsize\tabtight
\begin{tabular}{l r r r}
\toprule
Model & Gain (pp) & Frames & CI $\ni0$ \\
\midrule
TimeSformer & $-0.06$  & 6.13 & 8/8 \\
VideoMAE    & $-1.65$  & 6.70 & 6/8 \\
MC3-18      & $-1.92$  & 6.47 & 6/8 \\
R(2+1)D     & $-2.65$  & 6.62 & 4/8 \\
R3D-18      & $-2.68$  & 7.02 & 4/8 \\
ViViT       & $-4.96$  & 7.13 & 2/8 \\
SlowFast    & $-10.72$ & 7.13 & 1/8 \\
\bottomrule
\end{tabular}
\hspace{1.5em}
\begin{tabular}{l r r r l}
\toprule
TimeSformer & Top-1 & Frames & Gain (pp) & 95\% CI \\
\midrule
UCF-101       & 91.2 & 4.71 & $+0.14$ & $[-0.50,+0.50]$ \\
HMDB-51       & 80.4 & 5.21 & $+0.32$ & $[-1.32,+1.32]$ \\
Diving-48     & 37.9 & 5.73 & $+0.10$ & $[-2.13,+2.40]$ \\
DriveAct      & 67.7 & 5.84 & $+0.02$ & $[-3.07,+1.75]$ \\
EPIC-Kitchens & 31.8 & 6.43 & $-0.31$ & $[-1.56,+0.78]$ \\
AUTSL         & 66.9 & 6.64 & $+0.10$ & $[-0.81,+0.67]$ \\
SSv2          & 42.2 & 6.89 & $-0.03$ & $[-1.07,+0.83]$ \\
FineGym       & 65.0 & 7.62 & $-0.85$ & $[-2.35,+0.37]$ \\
\bottomrule
\end{tabular}
\end{table}

For TimeSformer the cascade is statistically indistinguishable from dense
inference on all 8 datasets at $4.7$--$7.6$ average frames
(Table~\ref{tab:sup_cascade}). Over all 56 pairs the mean change is $-3.5$pp at
$6.7$ frames, and 25 pairs have intervals that exclude zero. The losses
concentrate where stride-4 inference is itself degraded: FineGym (mean
$-13.8$pp) and SlowFast (mean $-10.7$pp, down to $-33.4$pp on FineGym). The
cascade is thus a useful operating point for backbones that are robust to
sparse sampling, in line with the temporal results of the main paper.

% ── S18: Convergence at low resolution ───────────────────────────────────────
\section{Convergence of Low-Resolution Fine-Tuning}
\label{sec:sup_convergence}

To check whether the 10-epoch budget limits the reported off-native accuracies,
we recovered the validation curve of each fine-tuning run from its training log
and asked whether the run was still improving when the budget ended.

\begin{table}[!htbp]
\centering
\caption{Convergence by target resolution. ``Still improving'' counts runs whose
best validation accuracy occurred at the final or penultimate epoch.}
\label{tab:sup_convergence}
\scriptsize\tabtight
\begin{tabular}{r r r r r}
\toprule
Resolution & Runs & Mean best epoch & Still improving & Gain in last 3 epochs \\
\midrule
48px  &  36 & 6.97 & 52.8\% & 1.38pp \\
96px  &  96 & 5.08 & 24.0\% & 4.21pp \\
112px &  93 & 4.94 & 16.1\% & 3.89pp \\
160px &  95 & 5.79 & 27.4\% & 3.27pp \\
224px & 183 & 5.66 & 23.0\% & 2.89pp \\
\bottomrule
\end{tabular}
\end{table}

At 48px over half the runs peak at the final epoch, against 16--27\% at the
other resolutions, so the budget binds most there. The accuracy still being
gained in the closing epochs at 48px is however the smallest of any resolution
($1.38$pp): these runs were converging slowly rather than being cut off
mid-ascent. The 48px results should be read with this limitation in mind.

% ── S19: TDS metric variants ─────────────────────────────────────────────────
\section{Sensitivity of the \texorpdfstring{\TDS{}}{TDS} Ranking to Its Definition}
\label{sec:sup_tds_variants}

\TDS{} uses the endpoints of the stride curve, an absolute accuracy drop, and a
positive part. Table~\ref{tab:sup_tds_variants} tests each choice with a variant
of the metric.

\begin{table}[!htbp]
\centering
\caption{Dataset demand under five definitions, and the rank correlation of each
against \TDS{} as defined in the main paper. \emph{Unclipped} removes
$[\cdot]^+$; \emph{AUC} integrates the whole stride curve trapezoidally in
$\log_2 s$; \emph{relative} divides by stride-1 accuracy; \emph{normalised}
divides by the headroom above chance.}
\label{tab:sup_tds_variants}
\scriptsize\tabtight
\begin{tabular}{l rrrrr}
\toprule
Dataset & \TDS{} & Unclipped & AUC & Relative & Normalised \\
\midrule
FineGym       & 55.92 & 55.92 & 24.33 & 76.92 & 78.03 \\
AUTSL         & 53.02 & 53.02 & 17.74 & 71.22 & 71.65 \\
SSv2          & 27.57 & 27.57 &  9.61 & 65.31 & 66.23 \\
DriveAct      & 21.88 & 21.88 &  8.84 & 31.87 & 33.35 \\
Diving-48     & 19.16 & 19.16 &  5.10 & 45.18 & 47.66 \\
HMDB-51       & 16.64 & 16.64 &  4.94 & 21.05 & 21.58 \\
EPIC-Kitchens &  9.74 &  9.74 &  3.02 & 26.86 & 27.73 \\
UCF-101       &  4.88 &  4.88 &  0.81 &  5.54 &  5.61 \\
\midrule
Spearman vs.\ \TDS{} & --- & 1.000 & 1.000 & 0.952 & 0.952 \\
\bottomrule
\end{tabular}
\end{table}

The positive part does not change any dataset score, so \emph{unclipped} is
numerically identical to \TDS{}. Integrating the full curve preserves the
ordering exactly. Normalising by baseline accuracy or by headroom above chance
exchanges only adjacent pairs (DriveAct with Diving-48, HMDB-51 with
EPIC-Kitchens); the three highest-demand datasets are identical under all five
definitions.

% ── S20: Multi-stride training ───────────────────────────────────────────────
\section{Multi-Stride Training as a Mitigation}
\label{sec:sup_tra}

We test whether exposure to sparse sampling during training mitigates the
stride-induced loss. Starting from the native checkpoint, we fine-tune with
$(\mathcal{C},s)$ drawn at random from the evaluation grid and compare against
the baseline checkpoint. Two augmented arms are trained. \textbf{Sampler} draws
its augmentation through the evaluation sampler of Sec.~\ref{sec:sup_padding},
including its end-repetition; \textbf{uniform} draws from the same grid but
resamples the available frames uniformly over the input. All arms are evaluated
with the evaluation sampler.

\begin{table}[!htbp]
\centering
\caption{Multi-stride training on AUTSL, coverage 100\%. ``Drop'' is the
stride-1 to stride-16 accuracy drop; ``recovered'' is the share of the baseline
drop regained at stride 16; ``dense cost'' is the accuracy given up at stride 1.}
\label{tab:sup_tra}
\scriptsize\tabtight
\begin{tabular}{l l rrrrr r r r}
\toprule
Model & Arm & $s{=}1$ & $s{=}2$ & $s{=}4$ & $s{=}8$ & $s{=}16$ & Drop & Recovered & Dense cost \\
\midrule
\multirow{3}{*}{TimeSformer}
 & baseline & 66.8 & 66.9 & 66.3 & 65.7 & 50.8 & 16.0 & --- & --- \\
 & sampler  & 63.5 & 63.6 & 63.1 & 62.6 & 55.3 &  8.2 & 28.4\% & $-3.3$ \\
 & uniform  & 63.1 & 63.4 & 63.4 & 62.3 & 52.6 & 10.4 & 11.7\% & $-3.7$ \\
\midrule
\multirow{3}{*}{R3D-18}
 & baseline & 75.0 & 75.3 & 68.0 & 27.2 &  6.5 & 68.5 & --- & --- \\
 & sampler  & 65.1 & 65.3 & 62.1 & 44.9 & 31.9 & 33.2 & 37.0\% & $-9.9$ \\
 & uniform  & 63.8 & 63.2 & 59.9 & 28.0 & 10.4 & 53.3 &  5.7\% & $-11.2$ \\
\bottomrule
\end{tabular}
\end{table}

Multi-stride training recovers $28.4\%$ of the stride-16 drop for TimeSformer
and $37.0\%$ for R3D-18, whose stride-16 accuracy rises from $6.5\%$ to
$31.9\%$ (Table~\ref{tab:sup_tra}). It is a trade rather than a full repair:
63--72\% of the drop remains, since at stride 16 only about four distinct frames
are available on AUTSL and training cannot restore evidence that is not
supplied, and it costs $3.3$pp (TimeSformer) and $9.9$pp (R3D-18) of
dense-sampling accuracy. Matching the training distribution to the deployment
sampler matters: at near-identical dense accuracy, the sampler arm recovers
$2.4\times$ (TimeSformer) and $6.5\times$ (R3D-18) more than the uniform arm.

{\small
\bibliographystyle{ieee}
\bibliography{main}
}

\end{document}
